\documentclass[11pt]{article}
\usepackage[utf8]{inputenc}
\usepackage[T1]{fontenc}
\usepackage{amsmath,amsfonts,amssymb}
\usepackage{graphicx}
\usepackage{hyperref}
\usepackage{cite}
\usepackage{geometry}
\geometry{margin=1in}

\title{Augmenting English with a Japanese-Inspired Agentive Noun System for Actor-Network Theory Modeling}

\author{Jun Kawasaki\\
\texttt{junkawasaki.com}}

\date{\today}

\begin{document}

\maketitle

\begin{abstract}
Actor–Network Theory (ANT) requires a consistent way to denote the myriad human and non-human actants in a network. However, English's current agent-noun formation system (using suffixes like -er, -or, -ist, etc.) is limited in productivity and semantic nuance. This paper proposes a theoretical and computational solution: an extension of English's morphological space inspired by Japanese agentive noun morphology (e.g. 〜者, 〜手, 〜家). We first discuss ANT's need for systematic actant labeling and analyze how English's existing agentive suffixes lack granularity and regularity. We then introduce a new set of generalized English agent-noun suffixes (e.g. -ôr, -eêr, -îst, -ânt, -êé, -hēad, -shōp, -êrē) modeled after Japanese, each denoting a distinct role type. An algorithmic word-formation model is presented to generate these agent nouns from arbitrary stems, with phonological rules for vowel linking, consonant assimilation, and exceptions. Use cases are explored in natural language processing pipelines (for automatically labeling semantic roles), knowledge graph ontologies (for naming relational nodes), and large language model (LLM) fine-tuning for relational semantics. We discuss limitations of this approach, implications for cross-linguistic morphology, and the potential of a more universal, high-resolution system of actant labeling. The proposed extension aims to enrich English lexicon systematically, enhancing the computational representation of complex networks of actors.
\end{abstract}

\section{Introduction}

In Actor–Network Theory (ANT), actants – the actors in a network – can be literally any entity, human or non-human, that is granted the source of an action. ANT's principle of generalized symmetry holds that all these entities should be described in the same terms, without privileging humans over nonhumans. In practice, this means when modeling or analyzing a sociotechnical network, we need a consistent linguistic scheme to label every actor or actant in the network. Whether the actant is a person, an organism, a machine, or an abstract artifact, it should be identifiable as playing a distinct role in the network's interactions.

However, the English language lacks a productive, systematic system of agentive nouns to uniformly label such actants. English does have agent nouns – words that mean "entity that does X" – typically formed by derivational suffixes (e.g. -er in driver from drive). Yet the system is irregular and limited. The suffix -er is the most common and can attach to many verbs (run → runner, build → builder), but it also overlaps with instruments (cutter, printer) and has exceptions. Other suffixes like -or, -ist, -ian, or Latinate -ant exist, but their usage is not fully productive or consistent. Many actions have no straightforward agentive noun (for example, steal has thief rather than stealer, teach has teacher but learn has no common learner for a student in everyday use, etc.), or they rely on periphrasis ("one who X"). Moreover, the semantic nuances of these suffixes are coarse – English does not routinely mark the difference between, say, a temporary performer of an action and a professional practitioner, except via context or additional descriptors.

This limitation is not just linguistic trivia; it poses a modeling challenge when we attempt to computationally represent networks of actors. In ANT-inspired data models or knowledge graphs, one might wish to label nodes by their roles ("driver", "regulator", "observer"). Using English's ad-hoc vocabulary can lead to inconsistent labels, polysemous terms, or gaps where no single-word label exists. A more systematic approach to naming actants could improve clarity and uniformity in such representations.

To find inspiration for such a system, we turn to Japanese. The Japanese language exhibits a rich set of agentive noun morphemes (often Sino-Japanese suffixes) that provide fine distinctions in meaning. For instance, Japanese can form agent nouns using 〜者 (-sha, denoting a person who does X in a general sense), 〜手 (-shu/te, denoting someone who does X as a skilled role or occupation), 〜家 (-ka, denoting an expert or professional in X, often artistic or academic), among others. These suffixes are highly productive within their categories and allow nuanced role identification.

In this paper, we propose to augment English with a similar multi-suffix system of agentive nouns. By introducing a set of novel derivational suffixes (marked with diacritics here to distinguish them from existing forms) and corresponding formation rules, we aim to emulate the productive precision of Japanese while fitting into English phonology and usage.

\section{Background}

\subsection{Actor–Network Theory and the Need for Consistent Actant Labeling}

Actor–Network Theory (ANT), developed by Bruno Latour, Michel Callon, John Law, and others, is a framework for describing heterogeneous networks of actors (actants) in both social and technical systems. In ANT, an actant is defined as anything that "acts or to which activity is granted by others". Crucially, an actant can be literally any entity – not only people, but also objects, technologies, organisms, or ideas.

This theoretical stance creates a practical linguistic challenge: How do we name or label these actants in a uniform way when we describe or model a network? In scholarly ANT literature, researchers often end up using descriptive phrases or borrowed terms to refer to various actants ("the spokesman", "inscription device", "ally", "intermediary", etc.). But when we move to computational representations – such as encoding an ANT analysis in a knowledge graph or building a simulation – we would ideally assign each actant a clear label indicating its role or action.

\subsection{Limitations of English Agent-Noun Formation}

English does have mechanisms to form agent nouns (nouns meaning "one who does X"), but these mechanisms are fragmentary and often non-systematic. The most common agentive suffix in English is -er, inherited from Old English (OE -ere) and productive to this day. We add -er to many verbs to indicate the person or thing that performs the action: driver (one who drives), builder (one who builds), scanner (something that scans). This suffix is relatively productive for a wide range of base verbs, especially Germanic-root verbs. However, -er has notable limitations:

\begin{itemize}
\item \textbf{Polysemy}: -er forms can denote not only persons but also instruments or machines (e.g. printer can mean a machine that prints, blender a device that blends). Context is needed to know if a cutter is a person using a knife or the tool itself. There is no morphological distinction in English between an agentive and an instrumental role – both are "the thing that does X".

\item \textbf{Limited -er scope}: -er typically attaches to verb stems, but it cannot attach to just any noun or adjective to mean "person associated with X" (that role is sometimes filled by other suffixes or compounding).

\item \textbf{Competing Suffixes}: English has several Latin-origin agentive suffixes that overlap in function with -er, such as -or, -ist, -ant, and -ian. These usually derive from Latin or French introductions and have more restricted usage.
\end{itemize}

The net effect of these limitations is that English lacks a single coherent paradigm for generating agentive nouns with fine-grained distinctions.

\subsection{High-Resolution Agentive Morphology in Japanese}

Japanese offers a compelling example of how a language can have multiple productive agentive noun forms to express subtle differences in role. Rather than a one-size-fits-all suffix, Japanese speakers choose from several suffixes (often written in kanji) depending on the nuance. Key agentive suffixes in Japanese include:

\begin{itemize}
\item \textbf{〜者 (-sha)} – General doer/actant: This suffix indicates an entity (usually a person, but not necessarily marked by profession) performing an action or being in a state. For example, 運転者 (untensha) means "driver" in the sense of "the person who is (currently) driving".

\item \textbf{〜手 (-shu/-te)} – Skill-based agent / professional: The character 手 means "hand", and by extension "person skilled with their hands in …" or someone who does something as a vocation. 運転手 (untenshu) is "driver" as an occupation (e.g. a taxi driver).

\item \textbf{〜家 (-ka)} – Expert or specialist, often artistic/academic: The character 家 (literally "house") in this usage means someone who occupies a certain realm or pursues something as a career or serious endeavor. For example, 作家 (sakka) means "writer/novelist" (one who writes books for a living).
\end{itemize}

This array of suffixes allows Japanese speakers to create specific agent nouns as needed, with each systematically derived and covering a spectrum of roles from temporary to persistent, amateur to professional.

\section{Proposed Extension: A System of Productive English Agentive Suffixes}

Inspired by the Japanese framework above, we propose a set of eight new agentive noun suffixes for English. Each suffix is designed to be highly productive and to carry a specific semantic nuance corresponding to a role type.

\begin{table}[h]
\centering
\begin{tabular}{|l|l|l|l|}
\hline
\textbf{Suffix} & \textbf{Japanese Analog} & \textbf{Intended Meaning} & \textbf{Example} \\
\hline
-ôr & $\sim$者 (-sha) & General doer/actant (contextual) & Xôr = entity that does X \\
-eêr & $\sim$手 (-shu/te) & Skilled doer or professional & Xeêr = practitioner of X \\
-îst & $\sim$家 (-ka) & Expert or devotee of X & Xîst = specialist of X \\
-ânt & (various) & Participant in X & Xânt = one involved in X \\
-êé & (cf. -ee) & Recipient of X & Xêé = one who is X-ed \\
-hēad & $\sim$長 (-chō) & Leader of X & Xhēad = leader of X-doers \\
-shōp & $\sim$屋 (-ya) & Dealer/provider of X & Xshōp = one who trades in X \\
-êrē & (cf. -ary) & Associated with X & Xêrē = functionary of X \\
\hline
\end{tabular}
\caption{Proposed Agentive Suffixes for English}
\end{table}

\subsection{Morphophonological Rules for Suffix Attachment}

To make these suffixes truly usable, we define how they attach to base words:

\begin{enumerate}
\item \textbf{Basic concatenation}: By default, form the agent noun by concatenating the base word and the suffix.

\item \textbf{Drop terminal "-e"}: If suffix begins with a vowel and base ends in silent "-e", drop the "e" (with exceptions for soft consonants).

\item \textbf{Consonant doubling}: For stressed monosyllables ending in vowel + consonant, double the consonant before vowel-initial suffixes.

\item \textbf{"-y to -i" conversion}: Change final "-y" to "-i" when preceded by a consonant and followed by a vowel-initial suffix.
\end{enumerate}

\section{Use Cases and Applications}

\subsection{Semantic Role Labeling and Information Extraction}

In NLP tasks like semantic role labeling, our system could generate descriptive labels on the fly. For example, if a sentence says "The robot analyzed the data and alerted the team," a system could label the robot as analyzeôr and alertôr, while the team becomes alertêé.

\subsection{Knowledge Graphs and Ontologies}

In designing ontologies, our system provides consistent naming options. Instead of generic "Sender" and "Receiver", one might have Messagôr and Messagêé, following a uniform pattern that makes the ontology easier to learn and extend.

\subsection{Large Language Model Fine-Tuning}

By fine-tuning LLMs on data using these forms, we can encourage models to internalize role distinctions. A model trained on "the teacheêr taught the teachêé" might better distinguish who initiates versus receives actions.

\section{Discussion and Limitations}

While theoretically powerful, this proposal faces several challenges:

\subsection{From an AI Perspective}
\begin{itemize}
\item \textbf{Vocabulary collision}: New forms may clash with existing words
\item \textbf{Processing cost}: Rule-based generation adds computational overhead
\item \textbf{Pre-trained model compatibility}: Novel words may not integrate well with existing embeddings
\end{itemize}

\subsection{From an ANT Perspective}
\begin{itemize}
\item \textbf{Oversimplification risk}: Forcing fluid roles into discrete categories
\item \textbf{Practical application difficulty}: Cumbersome for fieldwork
\item \textbf{Theoretical disconnection}: May seem separate from core ANT concepts
\end{itemize}

\subsection{From a Linguistics Perspective}
\begin{itemize}
\item \textbf{Naturalness concerns}: Neologisms may seem ungrammatical
\item \textbf{Over-differentiation}: Eight suffixes may be cognitively excessive
\item \textbf{Historical grounding}: Needs stronger connection to English evolution
\end{itemize}

\section{Conclusion}

English's current agentive noun system, while serviceable for everyday communication, falls short of providing systematic, productive, and semantically fine-grained actor labeling for complex networks. Our proposed Japanese-inspired suffix system offers a theoretical framework for addressing these limitations, with potential applications in NLP, knowledge representation, and cross-linguistic understanding.

While challenges exist regarding adoption and implementation, the systematic nature of this approach provides a foundation for more precise computational representation of actor roles in complex systems. Future work should focus on empirical validation through user studies and corpus analysis, as well as integration with existing NLP frameworks.

The ultimate goal is not to replace English's existing vocabulary, but to provide a systematic toolkit for generating precise role labels when needed, particularly in computational contexts where clarity and consistency are paramount.

\end{document} 